\documentclass[17pt,a4paper,landscape]{extarticle}
\usepackage[margin=2.35cm,top=0.12cm,left=1.05cm]{geometry}
\usepackage{amsmath}
\usepackage{amssymb}
\usepackage{tasks}
\usepackage{tikz}
\usepackage{xcolor}
\usepackage[sfdefault,lf]{FiraSans}
\renewcommand{\familydefault}{\sfdefault}
\setlength{\parindent}{0pt}
\pagestyle{empty}
\settasks{label=\textbf{\arabic*.}, label-width=2.2em, item-indent=2.95em, column-sep=2.1cm, after-item-skip=4.7em}
\renewcommand{\arraystretch}{1.15}
\everymath{\displaystyle}

\begin{document}
\sffamily
\boldmath

\boldmath

\boldmath

{\LARGE \textbf{Introduction to Circle Geometry --- Proficient}}
\begin{tasks}(3)
\task {Two angles stand on the same arc. One is \mbox{$70^\circ$} and the other is \mbox{$x+12^\circ$}. Find \mbox{$x$}.}
\task {The angle at the centre is \mbox{$4y^\circ$} and the angle at the circumference on the same arc is \mbox{$38^\circ$}. Find \mbox{$y$}.}
\task {A tangent meets a radius at an angle of \mbox{$a+18^\circ$}. Find \mbox{$a$}.}
\task {An angle in a semicircle is \mbox{$3b^\circ$}. Find \mbox{$b$}.}
\task {Two angles in the same segment are \mbox{$84^\circ$} and \mbox{$2x+6^\circ$}. Find \mbox{$x$}.}
\task {The angle at the centre is \mbox{$5m^\circ$} and the angle at the circumference on the same arc is \mbox{$55^\circ$}. Find \mbox{$m$}.}
\task {A tangent meets a radius at an angle of \mbox{$2p-6^\circ$}. Find \mbox{$p$}.}
\task {A circle has a chord as a diameter. Explain why any angle on the circumference standing on that diameter is a right angle.}
\task {Two angles stand on the same chord \mbox{$AB$} and are marked \mbox{$46^\circ$} and \mbox{$q$}. Find \mbox{$q$} and state the rule used.}
\task {The angle at the centre subtending arc \mbox{$AB$} is \mbox{$124^\circ$}. What is the angle at the circumference subtending the same arc?}
\task {The angle at the circumference subtending arc \mbox{$CD$} is \mbox{$33^\circ$}. What is the angle at the centre subtending the same arc?}
\task {A tangent touches a circle at \mbox{$T$}. The radius \mbox{$OT$} is drawn. A student says the angle between the tangent and \mbox{$OT$} could be \mbox{$88^\circ$}. Explain why this is impossible.}
\end{tasks}

\end{document}
